This study aims to analyze the relationship between peso-dollar exchange rate volatility and the trading volume of U.S. dollar futures contracts on the Mexican Derivatives Exchange (MexDer) over the 2009–2022 period. A quantitative, correlational-explanatory approach was used, with monthly series from Banco de México, MexDer and CME Group; exchange rate volatility was estimated with a GARCH(1,1) model, and its effect on trading volume was assessed through regressions with robust (Newey-West) standard errors for two subperiods. Results show that, between 2009 and 2015, a one-percentage-point increase in conditional volatility was associated with a 4.2\% rise in futures volume ($\beta = 0.042$, $p < 0.01$); however, between 2016 and 2022 the effect was no longer significant ($\beta = 0.011$, $p = 0.21$), while peso futures volume on the CME was negatively related to MexDer activity ($\beta = -0.351$, $p < 0.01$). Annual futures volume on MexDer fell by about 86\% over the period. The main limitation is the lack of monthly data on currency hedging arranged in the local over-the-counter (OTC) market. This research contributes to the literature on derivatives markets in emerging economies by showing that the contraction of MexDer is not driven by lower hedging demand, but by the migration of liquidity to foreign and OTC markets.

\hspace{1.27cm}\textit{Keywords:} derivatives market; dollar futures; exchange rate volatility; hedging; Mexico\\
\hspace*{1.27cm}\textit{JEL Classification:} G13, G15, F31
